% Part: first-order-logic
% Chapter: completeness
% Section: complete-sets

% Tam ve tutarlı kümelerin tanımı. Tamlık ispatında gereken özellikler,
% kullanılan ispat sistemine ilişkin bölümün provability.tex dosyasındadır.

\documentclass[../../../include/open-logic-section]{subfiles}

\begin{document}

\iftag{FOL}
      {\olfileid[tr]{fol}{com}{ccs}}
      {\olfileid[tr]{pl}{com}{ccs}}

\olsection{Tam ve Tutarlı \usetoken{P}{sentence} Kümeleri}

\begin{defn}[Tam küme]
\ollabel{def:complete-set} Bir !!{sentence}s kümesi~$\Gamma$,
\emph{!!{complete}}dır ancak ve ancak herhangi bir !!{sentence}~$!A$ için
ya $!A\in\Gamma$ ya da $\lnot !A\in\Gamma$ ise.
\end{defn}

\begin{explain}
!!^{complete} tümce kümeleri hiçbir soruyu yanıtsız bırakmaz. Her
!!{sentence}~$!A$ için $\Gamma$, $!A$'nın doğru mu yanlış mı olduğunu
``söyler.'' !!{complete} kümelerin önemi tamlık teoreminin ispatıyla
sınırlı değildir. Örneğin !!{complete} ve aksiyomlaştırılabilir bir
kuram her zaman karar verilebilirdir.
\end{explain}

\begin{explain}
!!^{complete} tutarlı kümeler tamlık ispatında önemlidir; çünkü her
tutarlı !!{sentence}s kümesi~$\Gamma$'nın !!a{complete} tutarlı
$\Gamma^*$ kümesinde bulunduğunu güvence altına alabiliriz.
!!^a{complete} tutarlı küme her !!{sentence}~$!A$ için ya $!A$'yı ya
da değillemesi~$\lnot !A$'yı içerir, fakat ikisini birden içermez. Bu,
özellikle \iftag{FOL}{atomik !!{sentence}s}{önerme değişkenleri} için
  doğrudur. Dolayısıyla !!a{complete} tutarlı bir kümeden
  \iftag{FOL}{---uygun biçimde !!{constant}s eklenerek genişletilmiş bir
  dilde---}{},
\iftag{FOL}{!!a{structure}}{!!a{valuation}} kurabiliriz. Bu
\iftag{FOL}{yapıdaki !!{predicate}s yorumunu}{değerlemenin önerme
değişkenlerine verdiği doğruluk değerlerini},~$\Gamma^*$ içinde bulunan
\iftag{FOL}{atomik !!{sentence}s}{!!{propositional variable}s}
belirler. Ardından bu \iftag{FOL}{!!{structure}}{!!{valuation}}
ifadesinin~$\Gamma^*$ içindeki bütün !!{sentence}s ifadelerini
(dolayısıyla~$\Gamma$ içindekilerin de tümünü) doğru kıldığı
gösterilebilir. Bu son olgunun ispatı için $\lnot !A\in\Gamma^*$ ancak
ve ancak $!A\notin\Gamma^*$; $(!A\lor !B)\in\Gamma^*$ ancak ve ancak
$!A\in\Gamma^*$ ya da $!B\in\Gamma^*$ olması vb. gerekir.
\end{explain}

Aşağıda $\Proves$ bağıntısının yansımalılık, monotonluk ve geçişlilik
özelliklerini çoğu kez örtük olarak kullanacağız (bkz.
\iftag{FOL}{%
  \tagrefs{prfSC/{fol:seq:ptn:sec},prfND/{fol:ntd:ptn:sec},prfAX/{fol:axd:ptn:sec},prfTab/{fol:tab:ptn:sec}}}{%
  \tagrefs{prfSC/{pl:seq:ptn:sec},prfND/{pl:ntd:ptn:sec},prfAX/{pl:axd:ptn:sec},prfTab/{pl:tab:ptn:sec}}}).

\begin{prop}
\ollabel{prop:ccs}
$\Gamma$'nın !!{complete} ve tutarlı olduğunu varsayalım. O zaman:
\begin{enumerate}
\item \ollabel{prop:ccs-prov-in} $\Gamma\Proves !A$ ise
  $!A\in\Gamma$.

\tagitem{prvAnd}{\ollabel{prop:ccs-and} $!A\land !B\in\Gamma$ ancak
  ve ancak hem $!A\in\Gamma$ hem de $!B\in\Gamma$ ise.}{}

\tagitem{prvOr}{\ollabel{prop:ccs-or} $!A\lor !B\in\Gamma$ ancak ve
  ancak $!A\in\Gamma$ ya da $!B\in\Gamma$ ise.}{}

\tagitem{prvIf}{\ollabel{prop:ccs-if} $!A\lif !B\in\Gamma$ ancak ve
  ancak $!A\notin\Gamma$ ya da $!B\in\Gamma$ ise.}{}
\end{enumerate}
\end{prop}

\begin{proof}
Aşağıdaki maddelerin tümünde~$\Gamma$'nın !!{complete} ve tutarlı
olduğunu varsayalım.
\begin{enumerate}
\item $\Gamma\Proves !A$ ise $!A\in\Gamma$.

$\Gamma\Proves !A$ olduğunu varsayalım. Aksine $!A\notin\Gamma$
olduğunu varsayalım. $\Gamma$ !!{complete} olduğundan
$\lnot !A\in\Gamma$ olur.
\iftag{FOL}{%
  \tagrefs{prfSC/{fol:seq:prv:prop:explicit-inc},
    prfND/{fol:ntd:prv:prop:explicit-inc},
    prfAX/{fol:axd:prv:prop:explicit-inc},
    prfTab/{fol:tab:prv:prop:explicit-inc}}}{%
  \tagrefs{prfSC/{pl:seq:prv:prop:explicit-inc},
    prfND/{pl:ntd:prv:prop:explicit-inc},
    prfAX/{pl:axd:prv:prop:explicit-inc},
    prfTab/{pl:tab:prv:prop:explicit-inc}}}
uyarınca~$\Gamma$ tutarsızdır. Bu,~$\Gamma$'nın tutarlı olduğu
varsayımıyla çelişir. Dolayısıyla $!A\notin\Gamma$ olamaz ve
$!A\in\Gamma$ olur.

\tagitem{defAnd}{}{%
\iftag{probAnd}{Alıştırma.}{%
$!A\land !B\in\Gamma$ ancak ve ancak hem $!A\in\Gamma$ hem de
$!B\in\Gamma$ ise:

İleri yön için $!A\land !B\in\Gamma$ olduğunu varsayalım.
\iftag{FOL}{%
  \tagrefs{prfSC/{fol:seq:ppr:prop:provability-land},
    prfND/{fol:ntd:ppr:prop:provability-land},
    prfAX/{fol:axd:ppr:prop:provability-land},
    prfTab/{fol:tab:ppr:prop:provability-land}}}{%
  \tagrefs{prfSC/{pl:seq:ppr:prop:provability-land},
    prfND/{pl:ntd:ppr:prop:provability-land},
    prfAX/{pl:axd:ppr:prop:provability-land},
    prfTab/{pl:tab:ppr:prop:provability-land}}}
sonucunun (1). maddesine göre $\Gamma\Proves !A$ ve
$\Gamma\Proves !B$ olur. \olref{prop:ccs-prov-in} uyarınca
$!A\in\Gamma$ ve $!B\in\Gamma$; istenen de buydu.

Ters yön için $!A\in\Gamma$ ve $!B\in\Gamma$ olsun.
\iftag{FOL}{%
  \tagrefs{prfSC/{fol:seq:ppr:prop:provability-land},
    prfND/{fol:ntd:ppr:prop:provability-land},
    prfAX/{fol:axd:ppr:prop:provability-land},
    prfTab/{fol:tab:ppr:prop:provability-land}}}{%
  \tagrefs{prfSC/{pl:seq:ppr:prop:provability-land},
    prfND/{pl:ntd:ppr:prop:provability-land},
    prfAX/{pl:axd:ppr:prop:provability-land},
    prfTab/{pl:tab:ppr:prop:provability-land}}}
sonucunun (2). maddesine göre $\Gamma\Proves !A\land !B$ olur.
\olref{prop:ccs-prov-in} uyarınca $!A\land !B\in\Gamma$.}}

\tagitem{defOr}{}{%
\iftag{probOr}{Alıştırma.}{%
Önce $!A\lor !B\in\Gamma$ ise $!A\in\Gamma$ ya da
$!B\in\Gamma$ olduğunu gösterelim. $!A\lor !B\in\Gamma$ iken
$!A\notin\Gamma$ ve $!B\notin\Gamma$ olduğunu varsayalım.
$\Gamma$ !!{complete} olduğundan $\lnot !A\in\Gamma$ ve
$\lnot !B\in\Gamma$ olur.
\iftag{FOL}{%
  \tagrefs{prfSC/{fol:seq:ppr:prop:provability-lor},
    prfND/{fol:ntd:ppr:prop:provability-lor},
    prfAX/{fol:axd:ppr:prop:provability-lor},
    prfTab/{fol:tab:ppr:prop:provability-lor}}}{%
  \tagrefs{prfSC/{pl:seq:ppr:prop:provability-lor},
    prfND/{pl:ntd:ppr:prop:provability-lor},
    prfAX/{pl:axd:ppr:prop:provability-lor},
    prfTab/{pl:tab:ppr:prop:provability-lor}}}
sonucunun (1). maddesine göre~$\Gamma$ tutarsız olur; bu bir
çelişkidir. Dolayısıyla $!A\in\Gamma$ ya da $!B\in\Gamma$.

Ters yön için $!A\in\Gamma$ ya da $!B\in\Gamma$ olduğunu varsayalım.
\iftag{FOL}{%
  \tagrefs{prfSC/{fol:seq:ppr:prop:provability-lor},
    prfND/{fol:ntd:ppr:prop:provability-lor},
    prfAX/{fol:axd:ppr:prop:provability-lor},
    prfTab/{fol:tab:ppr:prop:provability-lor}}}{%
  \tagrefs{prfSC/{pl:seq:ppr:prop:provability-lor},
    prfND/{pl:ntd:ppr:prop:provability-lor},
    prfAX/{pl:axd:ppr:prop:provability-lor},
    prfTab/{pl:tab:ppr:prop:provability-lor}}}
sonucunun (2). maddesine göre $\Gamma\Proves !A\lor !B$ olur.
\olref{prop:ccs-prov-in} uyarınca $!A\lor !B\in\Gamma$; istenen de
buydu.}}

\tagitem{defIf}{}{%
\iftag{probIf}{Alıştırma.}{%
İleri yön için $!A\lif !B\in\Gamma$ olduğunu, aksine
$!A\in\Gamma$ ve $!B\notin\Gamma$ olduğunu varsayalım. Bu varsayımlar
altında $!A\lif !B\in\Gamma$ ve $!A\in\Gamma$.
\iftag{FOL}{%
  \tagrefs{prfSC/{fol:seq:ppr:prop:provability-lif},
    prfND/{fol:ntd:ppr:prop:provability-lif},
    prfAX/{fol:axd:ppr:prop:provability-lif},
    prfTab/{fol:tab:ppr:prop:provability-lif}}}{%
  \tagrefs{prfSC/{pl:seq:ppr:prop:provability-lif},
    prfND/{pl:ntd:ppr:prop:provability-lif},
    prfAX/{pl:axd:ppr:prop:provability-lif},
    prfTab/{pl:tab:ppr:prop:provability-lif}}}
sonucunun (1). maddesine göre $\Gamma\Proves !B$. Fakat
\olref{prop:ccs-prov-in} uyarınca $!B\in\Gamma$ olur; bu da
$!B\notin\Gamma$ varsayımıyla çelişir.

Ters yön için önce $!A\notin\Gamma$ durumunu ele alalım. $\Gamma$
!!{complete} olduğundan $\lnot !A\in\Gamma$.
\iftag{FOL}{%
  \tagrefs{prfSC/{fol:seq:ppr:prop:provability-lif},
    prfND/{fol:ntd:ppr:prop:provability-lif},
    prfAX/{fol:axd:ppr:prop:provability-lif},
    prfTab/{fol:tab:ppr:prop:provability-lif}}}{%
  \tagrefs{prfSC/{pl:seq:ppr:prop:provability-lif},
    prfND/{pl:ntd:ppr:prop:provability-lif},
    prfAX/{pl:axd:ppr:prop:provability-lif},
    prfTab/{pl:tab:ppr:prop:provability-lif}}}
sonucunun (2). maddesine göre $\Gamma\Proves !A\lif !B$. Yine
\olref{prop:ccs-prov-in} uyarınca $!A\lif !B\in\Gamma$ elde ederiz.

Şimdi $!B\in\Gamma$ durumunu ele alalım.
\iftag{FOL}{%
  \tagrefs{prfSC/{fol:seq:ppr:prop:provability-lif},
    prfND/{fol:ntd:ppr:prop:provability-lif},
    prfTab/{fol:tab:ppr:prop:provability-lif},
    prfAX/{fol:axd:ppr:prop:provability-lif}}}{%
  \tagrefs{prfSC/{pl:seq:ppr:prop:provability-lif},
    prfND/{pl:ntd:ppr:prop:provability-lif},
    prfTab/{pl:tab:ppr:prop:provability-lif},
    prfAX/{pl:axd:ppr:prop:provability-lif}}}
sonucunun yine (2). maddesine göre $\Gamma\Proves !A\lif !B$.
\olref{prop:ccs-prov-in} uyarınca $!A\lif !B\in\Gamma$.}}
\end{enumerate}
\end{proof}

\tagprob{FOL}
\begin{prob}
\olref[fol][com][ccs]{prop:ccs} önermesinin ispatını tamamlayın.
\end{prob}
\tagendprob

\tagprob{notFOL}
\begin{prob}
\olref[pl][com][ccs]{prop:ccs} önermesinin ispatını tamamlayın.
\end{prob}
\tagendprob

\end{document}
